\section{Setting}\label{sec:setting}

Lattice observables are local functions of parallel transports; their limiting fields are distinguished by all separated correlations.

The faithful unitary realization $G\subset U(n_G)$ is fixed, with $G$ closed and its real Lie algebra $\g$ simple and nonabelian; connectedness is not assumed. Throughout,
\begin{equation}\label{eq:gauge-metric}
\langle X,Y\rangle_\g=-n_G^{-1}\operatorname{Re}\tr(XY),\qquad \dimG=\dim_\R\g\,.
\end{equation}
Here $\N=\{0,1,\ldots\}$, coordinate indices run from $1$ to $4$, $|x|$ is Euclidean and $|x|_\infty=\max_a|x^a|$. The block-spin penalty is a fixed $\blockpenalty>0$. Time is the fourth coordinate: $t(x)=x^4$, $\vec x=(x^1,x^2,x^3)$, $\minksquare(x)=(x^4)^2-|\vec x|^2$ and $\forwardcone=\{\mathbf p:|\vec{\mathbf p}|\le\mathbf p^4\}$. Fourier phases use $e^{i\mathbf p\cdot x}$ with the Euclidean pairing; Hilbert inner products are antilinear in the first variable. Spacetime integrals use Lebesgue measure. We include arity zero: empty products are $1$, and empty sums and the maximum of an empty degree list are $0$.

The ultraviolet, volume and renormalization indices are $N,p,m$, with $0\le m\le N$; blocking decreases $m$. Odd ratios $\ell_m\ge3$ give $D_m=\prod_{i<m}\ell_i$, $a_m=D_m^{-1}$ and $s_m=\log D_m$. The level-$m$ torus has $L_pD_m$ sites per side. Tokens $(d,\mathbf c)$ have insertions $O^{(d)}_{N,p}(\mathbf c,y)$; Wilson expectation is $\langle\cdot\rangle_{N,p}$. Fields $(d,A)$ are written $A$ when the degree is specified. Correlations $S(A_1(x_1),\ldots,A_n(x_n))$ and their derivatives concern separated points. Constants may depend on $G$, its realization, $\blockpenalty$ and the fixed construction; further dependence and uniformity are stated.
\paragraph{The lattice theory.} The matrix realization fixes the Wilson action and coupling normalization.

\begin{definition}[Gauge group and Casimir]\label{def:gauge-group}
A gauge group is an integer $n_G\ge1$, a closed subgroup $G\subset U(n_G)$ and a real subspace $\g\subset M_{n_G}(\C)$ such that
\begin{equation}\label{eq:lie-algebra}
X\in\g\quad\Longleftrightarrow\quad \exp(tX)\in G\quad\hbox{for every }t\in\R\,.
\end{equation}
Simplicity means that $[\g,\g]\ne0$ and every real subspace $I_{\rm id}\subset\g$ with $[\g,I_{\rm id}]\subset I_{\rm id}$ is $0$ or $\g$. With the fixed inner product and an orthonormal basis $(e_\alpha^{\g})_{\alpha=1}^{\dimG}$, define
\begin{equation}\label{eq:casimir}
C_{\g}=-\dimG^{-1}\sum_\alpha\operatorname{tr}_{\g}(\operatorname{ad}(e_\alpha^{\g})^2),\qquad b_0=\frac{11C_{\g}}{48\pi^2},\qquad \kappa_G=\frac{11C_{\g}}{24\pi^2}\,.
\end{equation}
Also put $\operatorname{pr}(u)=1-n_G^{-1}\operatorname{Re}\operatorname{tr}u$ and $\operatorname{ch}(u)=n_G^{-1/2}\|u-\mathrm{Id}_{n_G}\|_{\mathrm F}$, the chord to the identity, so $\operatorname{ch}(u)^2=2\operatorname{pr}(u)$.
\end{definition}

\begin{definition}[Configurations and tests]\label{def:configurations}
For $n\ge0$, let $\mathrm{Conf}_n=\{x\in(\R^4)^n:x_i\ne x_j\ (i\ne j)\}$ and $\mathrm{Gen}_n=\{x:x_i^a\ne x_j^a\ (i\ne j,\ a=1,\ldots,4)\}$. Write $\operatorname{sep}_\infty(x)=\min_{i\ne j}|x_i-x_j|_\infty$, with value $+\infty$ if $n\le1$. A separated test is a globally smooth complex function with compact closed support contained in $\mathrm{Conf}_n$. Write $\mathcal S_n=\mathcal S(\R^{4n};\C)$ and $\mathcal S=\mathcal S_1$. For a side $L>0$, put
\begin{equation}\label{eq:nohit}
\mathrm{NoHit}_{L,n}=\{x:x_i-x_j\notin L\Z^4\ (i\ne j)\}\,.
\end{equation}
The group $H_4\subset O(4)$ consists of the signed permutation matrices. Its reflection in the $a$th coordinate is $\mathsf s_a^{\rm refl}$.
\end{definition}

\begin{definition}[Wilson lattice theory]\label{def:wilson-theory}
For an integer $n_{\rm tor}\ge1$ put $\Lambda_{n_{\rm tor}}=(\Z/n_{\rm tor}\Z)^4$. A configuration assigns $U(y,a)\in G$ to each positive link $(y,a)$; a reversed link has inverse value. Define
\begin{align}\label{eq:wilson-law}
U_{y;ac}&=U(y,a)U(y+e_a,c)U(y+e_c,a)^{-1}U(y,c)^{-1},\\
E^{\rm W}_{n_{\rm tor}}(U)&=\sum_{y\in\Lambda_{n_{\rm tor}}}\sum_{a<c}\operatorname{pr}(U_{y;ac}),\\
\mu_{n_{\rm tor},\beta}(dU)&=Z_{n_{\rm tor},\beta}^{-1}e^{-\beta E^{\rm W}_{n_{\rm tor}}(U)}\prod_{y,a}dh(U(y,a))\,.
\end{align}
Here $dh$ is probability Haar measure on $G$ and $Z_{n_{\rm tor},\beta}$ is the integral of the numerator. A gauge transformation is $U^{\gamma_{\rm g}}(y,a)=\gamma_{\rm g}(y)U(y,a)\gamma_{\rm g}(y+e_a)^{-1}$. Lattice isometries act on links by transport, with inversion on reversed links. Let $r_{n_{\rm tor}}(j)$ be the representative of $j$ in $(-n_{\rm tor}/2,n_{\rm tor}/2]$.
\end{definition}

\begin{definition}[Wilson families and regularity]\label{def:wilson-family}
A Wilson family consists of strictly increasing positive integer sequences $(L_p),(D_N)$, real numbers $\beta_{N,p}$, probability Haar measure on $G$, a set $\mathcal T$ and functions $I_{N,p}(v,y;U)\in\C$ for $v\in\mathcal T$ and $y\in\Lambda_{L_pD_N}$. Here $(D_N)$ is arbitrary; for the constructed family $D_N=\prod_{i<N}\ell_i$. Put $a_N=D_N^{-1}$, $\hat y=a_N(r_{L_pD_N}(y^a))_{a=1}^4$ and $\mu_{N,p}=\mu_{L_pD_N,\beta_{N,p}}$.

The family is regular if:

\begin{description}[font=\normalfont,leftmargin=3.8em,labelindent=0pt,nosep]
\item[Reg1.] there is $b_*>0$ such that, for every $\varepsilon>0$, eventually in $N$, $|\beta_{N,p}/\log(1/a_N)-b_*|<\varepsilon$ for every $p$;
\item[Reg2.] every $U\mapsto I_{N,p}(v,y;U)$ is continuous;
\item[Reg3.] every such insertion is gauge invariant;
\item[Reg4.] for each $v$ there is a finite $\Sigma_v\subset\Z^4$, independent of $N,p,y$, such that the insertion depends only on links with both ends in $y+(\Sigma_v\bmod L_pD_N)$.
\end{description}
\end{definition}

\begin{definition}[Moments and ordered convergence]\label{def:ordered-convergence}
For $v\in\mathcal T^n$, set
\begin{equation}\label{eq:lattice-moments}
\latmoment_{N,p}(v;y)=\int\prod_iI_{N,p}(v_i,y_i;U)\,d\mu_{N,p}(U),\qquad \latsmear_{N,p}(v;f)=a_N^{4n}\sum_{y\in\Lambda_{L_pD_N}^n}f(\hat y)\latmoment_{N,p}(v;y)\,.
\end{equation}
A separated test $f$ is $p$-admissible if its support is contained in $((-L_p/2,L_p/2)^4)^n$. For $\mathsf v:\N\to\N$, ordered convergence to kernels $K_n(v;x)$ along $\mathsf v$ means that for every $n$, token list $v$ and separated test $f$: $f$ is $\mathsf v(k)$-admissible eventually in $k$, and there are numbers $C_p^\infty(v;f)$ such that $\latsmear_{N,p}(v;f)\to C_p^\infty(v;f)$ along the full sequence $N\to\infty$ at every $f$-admissible $p$, and $C_{\mathsf v(k)}^\infty(v;f)\to\int fK_n(v;\cdot)$. No condition is imposed on $C_p^\infty(v;f)$ at inadmissible volumes.
\end{definition}

\paragraph{Local observables and blocking.} Taylor normalization and transfer to a reference level fix the insertions.

\begin{definition}[Polynomial labels and local sections]\label{def:local-sections}
Let $\mathcal J$ be the space of $\g$-valued connection jets in radial gauge, written $\mathcal A_a(x)=\sum_{k\ge0}\mathcal A_{k,a}(x)$, where $\mathcal A_{k,a}$ is homogeneous of polynomial degree $k+1$ and $\sum_ax^a\mathcal A_{k,a}(x)=0$. Give the coordinates of $\mathcal A_k$ bidegree $(k+2,1)$: the first entry is engineering weight and the second counts polynomial factors. Let $V_d$ be the complex polynomials in these coordinates of engineering weight at most $d$ invariant under simultaneous $\operatorname{Ad}(G)$. Fix a basis homogeneous in both degrees, the supremum coefficient norm, conjugation and the induced $H_4$ action. Equivalently these are invariant polynomials in curvature jets, with curvature weight $2$ and covariant differentiation adding $1$.

A local section in degree $d$ is a linear map $\mathbf c\mapsto\localsection{d}[\mathbf c]$ to polynomials in the real and imaginary parts of link matrix entries on a fixed finite stencil $\Sigma_d\subset\Z^4$. Its placement $\localsection{d}[\mathbf c](y;U)$ evaluates it on the links of $y+\Sigma_d$. It is gauge invariant, commutes with conjugation, is $H_4$ covariant and sends the polynomial unit to $1$. Its Taylor normalization is the following equality for every radial jet $\mathsf j$ and $0\le k\le d$:
\begin{equation}\label{eq:local-taylor}
\left.\frac{d^k}{dh_{\rm lat}^k}\localsection{d}[\mathbf c](U(h_{\rm lat},\mathsf j))\right|_{h_{\rm lat}=0}
=\left.\frac{d^k}{dh_{\rm lat}^k}\mathbf c((h_{\rm lat}^{k'+2}\mathsf j_{k'})_{k'\ge0})\right|_{h_{\rm lat}=0}\,.
\end{equation}
Here $U(h_{\rm lat},\mathsf j)$ is parallel transport along the spacing-$h_{\rm lat}$ links of the radial polynomial potential through the required order, with transport equation $u'(t)=u(t)\mathcal A(t)$, $u(0)=1_G$.
\end{definition}

\begin{definition}[Constants and level data]\label{def:level-data}
The fixed numerical data are kernel parameters; stage data indexed by $j\ge4$ (a derivative order $\mathfrak q_j\ge6$, a spatial power $\varpi_j\ge2$, support exponent $\sigma_j=100$, path width $w_j^{\rm path}>0$, path bound $A_j^{\rm path}>0$, column bound $C_j\ge1$, drift allowance $\omega_j^{\rm RG}\ge0$ and positive clipping constants); odd stage ratios $\ell^{\mathrm{st}}_j$; exponents $\theta_j$ with $0<\theta_{J_m}\le J_m^{-1}$; entry bounds; a reference coupling $\bar\beta>1$; and entry levels and minimum periods. The orders $\mathfrak q_j$ and powers $\varpi_j$ tend to infinity. Section~\ref{sec:constants} specifies the kernel parameters, column bounds, minimum periods and order of choice. The requirements of Sections~\ref{sec:partII-B}--\ref{sec:partII-E}, \ref{sec:partII-G}, \ref{sec:partII-I} and~\ref{sec:partII-J} give lower bounds on stage ratios or entry levels using earlier data; raising entry levels preserves earlier requirements. Only the constants of the construction in Proposition~\ref{prop:U1}, with its densities, marked transfers and lattice estimates, are called admissible; no assertion is made for other choices.

A level system has $J_m\ge4$, $J_m\le J_{m+1}\le J_m+1$, $J_m\to\infty$, $\ell_m=\ell^{\mathrm{st}}_{J_{m+1}}$, and the spacings fixed above. The stage orders, powers and support exponents are nondecreasing along the levels; $\theta_j$ is antitone. Put
\begin{align}\label{eq:level-scales}
\hat\beta_m&=\bar\beta+\kappa_Gs_m,& \cutoffsize_m&=\exp((\log\log\hat\beta_m)^2),& t_m&=\hat\beta_m^{-1/2}\cutoffsize_m,\\
B_m&=\lceil(\log\hat\beta_m)^2\rceil,& \epsilon_m^{\rm reg}&=\hat\beta_m^{-11/10},& \epsilon_m^{\rm for}&=e^{-(\cutoffsize_m)^{1/5}},\\
\Lambda_m^{\mathrm{loop}}&=2(\cutoffsize_m)^{\theta_{J_m}},& \mathfrak B_m&=[\hat\beta_m/2,2\hat\beta_m]\,.
\end{align}
The entrance conditions are $\hat\beta_m>1$, $t_m\le1$, $L_pD_m\ge100(B_m+1)$, and $L_pD_{m+1}\ge100((J_{m+1}+3)\ell_m+1)$ for all $m,p$, together with
\begin{equation}\label{eq:entry-bounds}
\frac{\omega_{J_{m+1}}^{\rm RG}}{\log\ell_m}\le J_{m+1}^{-2},\qquad
\frac{\log(1+A_j^{\rm path}+(w_j^{\rm path})^{-1}+C_j+R_j^{\rm rel}+\sigma_j)}{\log\ell_m}\le j^{-2},\quad j=J_{m+1}\,,
\end{equation}
where $R_j^{\rm rel}$ is the clipping relaxation length. The stage ratios are nondecreasing. The conditions $\log\ell_m\to\infty$ and $\log\ell_m/s_m\to0$ are required only as limits; the quotient at $m=0$ plays no role.
\end{definition}

\begin{definition}[Block-spin transformation]\label{def:block-spin}
For a step of the fixed level system, $U,V$ are link configurations on $\Lambda_{L_pD_{m+1}},\Lambda_{L_pD_m}$. Let $\alpha_{\rm tag}$ range over the finite frame and predictor tags of Section~\ref{sec:partII-B}, $\lambda^{\rm tag}(U,\alpha_{\rm tag})$ be their product probability mass and $Q_e^{\rm pred}(U,\alpha_{\rm tag})\in G$ the predictor on coarse link $e$. The cutoffs ensure unique minimizing branches and nonempty eligible axial tags. These choices are fixed for the level system. Define
\begin{align}\label{eq:block-kernel}
\blockkernel{m}_\beta(U,V)&=\sum_{\alpha_{\rm tag}}\lambda^{\rm tag}(U,\alpha_{\rm tag})\prod_{e\text{ coarse}}\frac{e^{-\blockpenalty\beta\wilsonterm(V_eQ_e^{\rm pred}(U,\alpha_{\rm tag})^{-1})}}{\int_Ge^{-\blockpenalty\beta\wilsonterm(u)}\,dh(u)},\\
(\blockmap{m}_\beta f)(V)&=\int\blockkernel{m}_\beta(U,V)f(U)\,dU\,.
\end{align}
Here $dU,dV$ are product probability Haar measures. Since the tag masses sum to one and $dh$ is invariant, $\int\blockkernel{m}_\beta(U,V)\,dV=1$.
\end{definition}

\begin{definition}[Representations and their admissible class]\label{def:representations}
Distances in this definition use the unit-lattice sup metric. At level $m$, flag every site incident to a plaquette whose chord exceeds $t_m$, and let $\flags_m(U)$ be the set of flagged sites. A support $X$ is healthy when it contains no flagged site. A regular activity has a support $X_i$, anchor $o_i\in X_i$ and a Borel gauge-invariant real function $E_i^{\rm reg}$, local on $X_i$, smooth through order $\mathfrak q_{J_m}$ in exponential left-link increments on healthy configurations, with the corresponding ordered derivatives existing there. It vanishes at the trivial configuration. A support is nonwrapping at $o_i$ if $10\operatorname{dist}(x,o_i)$ is less than the torus side for every $x\in X_i$. On such a support the Hessian of $\mathcal A\mapsto E_i^{\rm reg}(\exp\mathcal A)$ at zero vanishes on every pair of affine directions $\dot{\mathcal A}_{(y,a)}=u_a+\sum_c r_{L_pD_m}(y-o_i)^cv_{ac}$, using the centered lift $r_{L_pD_m}$. A forced activity has a support $Y_i$, nonempty flag set $I_i\subset Y_i$ and a Borel gauge-invariant function $K_i^{\rm for}$, vanishing unless $Y_i\cap\flags_m(U)=I_i$ and a fixed Borel gauge-invariant test holds. Both $K_i^{\rm for}$ and that test take equal values on configurations agreeing on the links inside $Y_i$ and the flags in $Y_i$. Both activity families are countable and equivariant under translations and $H_4$, by permutations of their indices, including supports, anchors, flag sets, tests and values.

For a support $X$, let $\operatorname{cov}_{B_m}(X)$ be the minimum number of radius-$B_m$ lattice balls with connected union covering $X$. Let $\|E_i^{\rm reg}\|_m$ be the supremum of all ordered left derivatives through order $\mathfrak q_{J_m}$, over healthy configurations and directions satisfying $|\dot{\mathcal A}_e|\le t_m(1+\operatorname{dist}(e^-,o_i))^{\varpi_{J_m}}$. The activity bounds mean the existence of nonnegative majorants $e_i^{\rm maj}\ge\|E_i^{\rm reg}\|_m$, $k_i^{\rm maj}\ge\sup_U|K_i^{\rm for}(U)|$ such that, for every site $x$, the following series converge and satisfy
\begin{equation}\label{eq:activity-bounds}
\sum_{o_i=x}e^{\sigma_{J_m}\operatorname{cov}_{B_m}(X_i)}e_i^{\rm maj}\le\epsilon_m^{\rm reg},\qquad
\sum_{x\in Y_i}e^{\sigma_{J_m}\operatorname{cov}_{B_m}(Y_i)}k_i^{\rm maj}\le\epsilon_m^{\rm for}\,.
\end{equation}
A representation $\rho=(\beta,v_{\rm vac},\pathaction,E^{\rm reg},K^{\rm for})$ has the fixed path action $\pathaction=\pathaction_{N,m}$ and density
\begin{multline}\label{eq:represented-density}
\density\rho(U)=\exp\left\{v_{\rm vac}|\Lambda_{L_pD_m}|-\pathaction(\beta,U)-\sum_i\mathbf1_{X_i\text{ healthy}}E_i^{\rm reg}(U)\right\}\\
{}\cdot\sum_{\substack{\mathcal A_{\rm for}\subset\N\text{ finite}\\ Y_i\ (i\in\mathcal A_{\rm for})\text{ pairwise disjoint}\\ \bigcup_{i\in\mathcal A_{\rm for}}I_i=\flags_m(U)}}\prod_{i\in\mathcal A_{\rm for}}K_i^{\rm for}(U)\,.
\end{multline}
The path action and its arrays, with $c_{\rm dir}>16c_{\rm lin}>0$, are defined in Sections~\ref{sec:partII-A}--\ref{sec:partII-D}, independently of $p$. A slot is a finite lattice path followed by a plaquette, evaluated as $\operatorname{proj}_{\g}(U_\square-U_\square^{-1})/2$ transported to the path origin. For these slots $\mathcal F_{\mathfrak p i}$, rows $\mathcal L_{\mathfrak p}^{\rm path}=\sum_i r_{\mathfrak p i}\mathcal F_{\mathfrak p i}$ and invariant $k$-linear atoms $\mathfrak a_i^{(k)}$, the bounds are
\begin{equation}\label{eq:path-bounds}
\sum_i|r_{\mathfrak p i}|e^{w_{J_m}^{\rm path}\operatorname{length}(\mathcal F_{\mathfrak p i})}\le A_{J_m}^{\rm path},\qquad
\sum_i|\mathfrak a_i^{(k)}|e^{w_{J_m}^{\rm path}\operatorname{length}(\mathfrak a_i^{(k)})}\le
\begin{cases}\Lambda_m^{\rm loop},&k=2,\\ A_{J_m}^{\rm path},&k=3,4.\end{cases}
\end{equation}
Slot length is path length plus four; atom length is the sum of its slot lengths, and atom magnitude is the sum of absolute coefficients in the fixed colour basis. If $s_{\mathfrak p\mathfrak q}$ sums the row coefficients of orientation $\mathfrak q$, kinetic normalization means $c_{\rm dir}\delta_{\mathfrak q\mathfrak r}+\sum_{\mathfrak p}s_{\mathfrak p\mathfrak q}s_{\mathfrak p\mathfrak r}=\delta_{\mathfrak q\mathfrak r}$. For each ordered pair of orientations and colour components, the sum of quadratic-atom coefficients and its first moments in the difference of slot offsets vanish, with absolute convergence.

An admissibility datum consists of subsets $\mathcal A_{N,m,p}$ of these representations, initial subsets $\mathcal I_{N,p}\subset\mathcal A_{N,N,p}$ and relations $\mathcal R_{N,m,p}\subset\mathcal A_{N,m,p}\times\mathcal A_{N,m+1,p}$. A representation is admissible when it belongs to $\mathcal A_{N,m,p}$. A trajectory is a tuple $(\rho_m)_{0\le m\le N}$ with $\rho_N\in\mathcal I_{N,p}$ and $(\rho_m,\rho_{m+1})\in\mathcal R_{N,m,p}$ for every $m<N$. Proposition~\ref{prop:U1} supplies these subsets and relations, with its continuation property on the chosen subsets; the activity bounds alone do not define membership.

Along a curve the index sets, supports, anchors and flag sets are fixed. Its $\beta,v_{\rm vac}$ are continuous; differences of regular derivatives and forced functions admit majorants whose sums with the support weights of~\eqref{eq:activity-bounds} tend to zero as the parameter approaches any point. No continuity across flag boundaries is required.
\end{definition}

\begin{definition}[Renormalized insertions]\label{def:insertions}
For each $d$ choose a reference level $\reflevel{d}$ with $J_{\reflevel{d}}\ge d+3$ and $\mathfrak q_{J_m}\ge d+1$, $\varpi_{J_m}\ge d$ for every $m\ge \reflevel{d}$. Let $\factorscale_t$ multiply a basis vector of bidegree $(w,k)$ by $t^k$, counting polynomial factors. For $N\ge\reflevel{d}$, use the trajectory $\rho_{N,p,m}$ and couplings $\beta_{N,p,m}$ of Proposition~\ref{prop:U2}. The marked densities, linear in terminal input $\mathbf v\in V_d$, are defined at the origin by
\begin{align}\label{eq:marked-transfer}
H^{\rm ins}_{N,p,N;d}(\mathbf v;U)&=\density\rho_{N,p,N}(U)\,\localsection{d}[\factorscale_{t_N^{-1}}\mathbf v](0;U),\\
H^{\rm ins}_{N,p,m;d}(\mathbf v;\cdot)&=\ell_m^d\blockmap{m}_{\beta_{N,p,m+1}}H^{\rm ins}_{N,p,m+1;d}(\mathbf v;\cdot)\,.
\end{align}
The terminal factors $a_N^d$ and $a_N^{-d}$ cancel. The coefficient at level $\reflevel{d}$ is $\mathsf Z_{N,p,d}\mathbf v$, defining an invertible map $\mathsf Z_{N,p,d}:V_d\to V_d$. At $N<\reflevel{d}$ choose any invertible map. For every cutoff and site $y$ set
\begin{equation}\label{eq:renormalized-insertion}
O^{(d)}_{N,p}(\mathbf c,y;U)=a_N^{-d}\localsection{d}[\factorscale_{t_N^{-1}}\mathsf Z_{N,p,d}^{-1}\mathbf c](y;U)\,.
\end{equation}
The distinguished unit token $\mathbf u_d$ is $a_{\reflevel{d}}^d$ times the polynomial unit in coefficient coordinates; its insertion is $1$ for $N\ge\reflevel{d}$. The quotient unit of Definition~\ref{def:euclidean-table} is $[\mathbf u_d]$. The marked-density remainder norms are given in Section~\ref{sec:partII-E}.
\end{definition}

\begin{definition}[Chessboard roots and differences]\label{def:chessboard}
For a tiling $\mathcal C^{\rm tile}$ of the fine torus into congruent cells, reflect a base-cell observable $F_{\rm cell}$ into every cell, applying conjugation on odd reflections, and denote these copies by $(F_{\rm cell})_C$. Its chessboard root is
\begin{equation}\label{eq:chessboard-root}
\|F_{\rm cell}\|_{\mathcal C^{\rm tile}}=\left|\left\langle\prod_{C\in\mathcal C^{\rm tile}}(F_{\rm cell})_C\right\rangle_{N,p}\right|^{1/|\mathcal C^{\rm tile}|}\,.
\end{equation}
Interior means dependence only on links with both ends strictly inside the base cell. The centered tiling $\mathcal C_m^{\rm tile}(y)$ has physical cell side $a_m$ and base cell centered at $y$. For a list $\boldsymbol l=((i_1,a_1),\ldots,(i_k,a_k))$ of slot-axis pairs put $\Delta_{\boldsymbol l}=\prod_{\alpha=1}^k a_N^{-1}(\tau_{i_\alpha,a_\alpha}-1)$, where $\tau_{i,a}$ shifts lattice slot $i$ forward by one link. Let $\partial_{\boldsymbol l}$ be the corresponding continuum derivative. Physical torus distance is $\min_{z\in L_p\Z^4}|\hat y-\hat y'+z|_\infty$.
\end{definition}

\paragraph{Euclidean fields and stress.} The field quotient uses all spectators; shell fluxes normalize stress.

\begin{definition}[Euclidean tables and field spaces]\label{def:euclidean-table}
A table $S$ assigns a complex multilinear function of $\mathbf c_i\in V_{d_i}$ to each list of degrees and points. A word is a finite list $w=((x_1,\mathbf c_1),\ldots,(x_n,\mathbf c_n))$ with distinct positions, and $S(w)=S(\mathbf c_1(x_1),\ldots,\mathbf c_n(x_n))$. Write $ww'$ for concatenation, $hw=((hx_i,h\mathbf c_i))_i$ for $h\in H_4$, and $w^*=((x_n,\overline{\mathbf c_n}),\ldots,(x_1,\overline{\mathbf c_1}))$. A word is positive in axis $a$ if every $x_i^a>0$. The Euclidean properties of $S$ are:

\begin{description}[font=\normalfont,leftmargin=3.8em,labelindent=0pt,nosep]
\item[E1.] simultaneous permutation symmetry, $S_0=1$, and deletion of a distinguished unit slot $\mathbf u_d$.
\item[E2.] $S(h\mathbf c;hx)=S(\mathbf c;x)$ for $h\in H_4$, and $S(\overline{\mathbf c};x)=\overline{S(\mathbf c;x)}$.
\item[E3.] $S(\mathbf c;x+z)=S(\mathbf c;x)$ for every $z\in\R^4$.
\item[E4.] for positive words $w_i$ and coefficients $a_i^{\rm RP}$, $\sum_{i,j}\overline{a_i^{\rm RP}}a_j^{\rm RP}S((\mathsf s_a^{\rm refl}w_i^*)\,w_j)\ge0$ for each $a$.
\item[E5.] the same inequality holds for finite sums of words smeared against joint smooth compactly supported tests on positive-side separated configurations. Every cross-term integral is absolutely convergent.
\item[E6.] every kernel is smooth on $\mathrm{Conf}_n$.
\item[E7.] for every $n,d_i,\mathbf c_i$ and positional derivative order $k$, there are $C\ge0$ and an integer $q_{\rm grow}\ge0$ such that $\|D^kS(\mathbf c;x)\|\le Cr^{-q_{\rm grow}}$ whenever $0<r\le1$ and $\operatorname{sep}_\infty(x)\ge r$.
\end{description}
For such a table, with $n\ge0$ and $\boldsymbol v$ any list of $n$ spectator tokens, define
\begin{equation}\label{eq:field-quotient}
\begin{split}
\nullspace{d}(S)&=\{\mathbf c\in V_d:S(\mathbf c(x_0),\boldsymbol v(x))=0\text{ for every }n,\boldsymbol v\text{ and }(x_0,x)\in\Conf{n+1}\},\\
F_d&=V_d/\nullspace{d}(S)\,.
\end{split}
\end{equation}
Equip each finite-dimensional quotient with a norm. The response of $[\mathbf c]$ is its function of all spectators and separated configurations in~\eqref{eq:field-quotient}; their linear span is the response label space $\mathcal L$. Inclusions $\jmath_{d,e}$, derivatives $\partial_a$, conjugation and units are characterized by equality, differentiation, conjugation and deletion of slots in all correlations. Proposition~\ref{prop:P5} proves their existence and the directed-union property of $\mathcal L$. The quotient unit is $\unitfield{d}=[\mathbf u_d]$.
\end{definition}

\begin{definition}[Free Maxwell fields and frames]\label{def:maxwell}
Let $\dimG$ independent real Euclidean Maxwell curvatures have potential covariance $\delta_{\alpha\beta}\delta_{ac}\Gamma(x-y)$, $\Gamma(x)=1/(4\pi^2|x|^2)$. Their curvature covariance is obtained by applying $\partial_a\mathcal A_c-\partial_c\mathcal A_a$ in each argument. The on-shell curvature jets satisfy antisymmetry, the linear Bianchi equations and $\sum_a\partial_aF_{ac}=0$. Let $\maxwellspace{d}$ be the space of $\operatorname{Ad}(G)$-invariant complex Wick polynomials in these jets of weight at most $d$, modulo these equations; $S^{\mathrm{Max}}$ is their separated Wick table. This defines its units, conjugation, inclusions, derivatives and $O(4)$ actions $R_d^{\mathrm{Max}}$. Put $\curvaturesquare=\frac14\sum_{a,c,\alpha}:F^\alpha_{ac}F^\alpha_{ac}:$. Let $T^{\mathrm{Max}}_{ab}\in\maxwellspace{4}$ be the multiple of $T^{\rm can}_{ab}=\sum_{c,\alpha}:F^\alpha_{ac}F^\alpha_{bc}:-\delta_{ab}\curvaturesquare$ fixed by the translation shell law of Definition~\ref{def:stress-laws} with $\mathsf p_b=\partial_b$.

A scaling sequence consists of positive radii $r_k\to0$ and complex linear isomorphisms $\scaleframe{k}{d}:\maxwellspace{d}\to F_d$. It is stable if, at each fixed degree horizon, eventually these maps preserve conjugation, $H_4$, units and inclusions exactly. Scaling dimensions are included in $\scaleframe{k}{d}$; an additional factor $r_k^{d}$ is not put on ordinary comparison slots.
\end{definition}

\begin{definition}[Words and slab operators]\label{def:os-spaces}
For a table $S$ with E1--E4, quotient the complex span of positive words by the null space of $\langle w,w'\rangle=S((\mathsf s_a^{\rm refl}w^*)\,w')$ and complete, as in~\cite[Section~4.1, p.~91]{OsterwalderSchrader1973}. This is $\mathcal H^{\mathrm{OS}}_a$. The empty word is $\Omega_a^{\rm OS}$. Positive translation in direction $a$ induces the contraction semigroup $e^{-t_{\rm heat}H_a^{\rm OS}}$; put $\mathcal H^0_a=\ker H_a^{\rm OS}$ and let $P^0_a$ be its orthogonal projection.

A slab word of width $h_{\rm sl}>0$ has $0<x_i^4<h_{\rm sl}$. It has buffer $r$ if $\operatorname{sep}_\infty(x)\ge r$ and $r\le x_i^4\le h_{\rm sl}-r$. For such a word $w$, let $A_ww'=w\,\tau_{h_{\rm sl} e_4}w'$ on positive words, where $\tau_z$ translates every position by $z$, and let $w^\dagger=\tau_{h_{\rm sl} e_4}\mathsf s_4^{\rm refl}w^*$. The repeated reflected word with exterior $w_{\rm ext}$ is $(A_{w^\dagger}A_w)^{n_{\rm rep}}w_{\rm ext}$. A slab operator is the bounded extension $\mathsf B_{h_{\rm sl}}(w)$ of $A_w$, when it exists. Its compression is $M^0(w)=P^0_4\mathsf B_{h_{\rm sl}}(w)|_{\mathcal H^0_4}$, after any translation placing the word in a slab. Proposition~\ref{prop:P8} proves existence and independence.
\end{definition}

\begin{definition}[Shell fluxes and stress laws]\label{def:stress-laws}
For separated $x_0,x_1,\ldots,x_n$, let $\chi\in C_c^\infty(\R^4;\R)$ equal $1$ near $x_0$, with $x_i\notin\operatorname{supp}\chi$ for $i>0$. For a tensor $J_{ab}\in F_4$, $A\in F_d$, spectator list $\boldsymbol B$, and $K_{\rm rot}\in\mathfrak{so}(4)$, define
\begin{align}\label{eq:shell-flux}
\mathcal P_b^J(\chi;A,\boldsymbol B;x)&=-\sum_a\int\partial_a\chi(y)S(J_{ab}(y),A(x_0),\boldsymbol B(x))\,dy,\\
\mathcal Y_{K_{\rm rot}}^J(\chi;A,\boldsymbol B;x)&=-\sum_{a,b}\int\partial_a\chi(y)(K_{\rm rot}(y-x_0))_bS(J_{ab}(y),A(x_0),\boldsymbol B(x))\,dy\,.
\end{align}
A conserved tensor is real, symmetric and satisfies $\sum_a\partial_aJ_{ab}=0$ in $F_5$. It is centered when $S(J_{ab}(x))=0$ for all $a,b,x$. It is $H_4$ tensorial when $hJ_{ab}=\sum_{i,j}h_{ia}h_{jb}J_{ij}$. Raw shell laws for $\mathsf p_b,\mathsf Y_{K_{\rm rot}}$ mean $\mathcal P_b^J=S((\mathsf p_bA)(x_0),\boldsymbol B(x))$ and $\mathcal Y_{K_{\rm rot}}^J=S((\mathsf Y_{K_{\rm rot}}A)(x_0),\boldsymbol B(x))$ for every such choice. For an $O(4)$ representation $R_d$, write $\mathsf L_d(K_{\rm rot})=\left.\frac d{dt}R_d(e^{tK_{\rm rot}})\right|_{t=0}$. The normalized stress laws use $\mathsf p_b=\partial_b$ and $\mathsf Y_{K_{\rm rot}}=\mathsf L_d(K_{\rm rot})$. These integrals avoid all insertion points.
\end{definition}

\paragraph{Relativistic reconstruction.} Phase correlations determine relativistic fields.

\begin{definition}[Relativistic data]\label{def:relativistic-data}
A relativistic datum is a complex Hilbert space $\mathcal H$, a complex vector space $\mathcal L$ (the response space of Definition~\ref{def:euclidean-table} in Theorem~\ref{thm:B}), a complex linear subspace $\mathcal D\subset\mathcal H$, $\Omega\in\mathcal D$, fields $\Phi(\ell_{\rm lab},f):\mathcal D\to\mathcal D$ complex linear in both $\ell_{\rm lab}\in\mathcal L$ and $f\in\mathcal S$, an antilinear label map $\ell_{\rm lab}\mapsto\ell_{\rm lab}^*$, and strongly continuous unitary translations $\mathsf T(z_{\rm tr})$. It includes a Borel projection-valued measure $\mathsf E$ on $\R^4$, extended by zero to non-Borel subsets, with diagonal finite measures $\mu_\psi(A)=\langle\psi,\mathsf E(A)\psi\rangle$ for Borel $A$.

It also includes a topological group $\mathcal P$ identified by a homeomorphism $\kappa_{\rm P}$ with $SO^+(1,3)\ltimes\R^4$, respecting identity and multiplication; write $\kappa_{\rm P}(\gamma)=(\Lambda_\gamma,z_\gamma)$. Here $\Lambda$ preserves $\mathsf m$, has determinant $1$, and maps $\overline V_+$ onto itself. There are strongly continuous unitaries $\mathsf U(\gamma)$ forming a representation, with $\mathsf U(1,z_{\rm tr})=\mathsf T(z_{\rm tr})$, and linear bijections $\rho_{\rm lab}(\gamma)$ of $\mathcal L$. Put $(\gamma\cdot f)(x)=f(\Lambda_\gamma^{-1}(x-z_\gamma))$. No continuity, finite-dimensionality, translation triviality or involution property is imposed on the label data in this definition.
\end{definition}

Only Borel sets enter the spectral axioms below; assigning zero elsewhere fixes a convention for the same Borel projection-valued measure.

\begin{definition}[Wightman axioms]\label{def:wightman}
The datum in Definition~\ref{def:relativistic-data} satisfies the following axioms.

\begin{description}[font=\normalfont,leftmargin=3.8em,labelindent=0pt,nosep]
\item[W0a.] $\langle\psi,\mathsf T(z_{\rm tr})\psi\rangle=\int e^{i\mathbf p\cdot z_{\rm tr}}\,d\mu_\psi(\mathbf p)$ for every $\psi,z_{\rm tr}$.
\item[W0b.] $\mathsf E(\R^4\setminus\overline V_+)=0$.
\item[W0c.] every vector fixed by all translations belongs to $\C\Omega$.
\item[W1a–c.] $\|\Omega\|=1$, $\mathcal D$ is dense, and $\mathcal D$ equals the complex span of all $\Phi(\ell_{\rm lab,1},f_1)\cdots\Phi(\ell_{\rm lab,k},f_k)\Omega$, including $k=0$.
\item[W1d.] for each $\ell_{\rm lab}$ and $\psi\in\mathcal D$, $f\mapsto\Phi(\ell_{\rm lab},f)\psi$ is continuous from $\mathcal S$ to $\mathcal H$ in norm.
\item[W1e.] $\langle\Phi(\ell_{\rm lab},f)\psi,\psi'\rangle=\langle\psi,\Phi(\ell_{\rm lab}^*,\bar f)\psi'\rangle$ on $\mathcal D$.
\item[W2a.] $\rho_{\rm lab}(1)=1$ and $\rho_{\rm lab}(\gamma\gamma')=\rho_{\rm lab}(\gamma)\rho_{\rm lab}(\gamma')$.
\item[W2b–d.] $\mathsf U(\gamma)\mathcal D\subset\mathcal D$, $\mathsf U(\gamma)\Omega=\Omega$, and $\mathsf U(\gamma)\Phi(\ell_{\rm lab},f)\psi=\Phi(\rho_{\rm lab}(\gamma)\ell_{\rm lab},\gamma\cdot f)\mathsf U(\gamma)\psi$.
\item[W3.] if every $x\in\operatorname{supp}f$ and $y\in\operatorname{supp}f'$ have $\mathsf m(x-y)<0$, then $[\Phi(\ell_{\rm lab},f),\Phi(\ell'_{\rm lab},f')]\psi=0$ for every $\ell_{\rm lab},\ell'_{\rm lab}$ and $\psi\in\mathcal D$.
\end{description}
\end{definition}

\begin{definition}[Hierarchy data]\label{def:hierarchy}
A hierarchy over a Wilson family is the following data, with no further conditions implicit:

\begin{description}[font=\normalfont,leftmargin=3.8em,labelindent=0pt,nosep]
\item[H1.] a relativistic datum as in Definition~\ref{def:relativistic-data};
\item[H2.] finite-dimensional complex normed spaces $F_d$ and linear maps $\iota_d:F_d\to\mathcal L$; put $\mathbb F=\bigsqcup_{d\ge0}F_d$, $\deg(d,A)=d$, $\iota(d,A)=\iota_dA$;
\item[H3.] a map $\pi_{\rm quo}:\mathcal T\to\mathbb F$;
\item[H4.] a compact metric space $\phases$, a Borel measure $\nu$ and a point $\xi_*\in\phases$;
\item[H5.] functions $S_n^\zeta:\mathbb F^n\times(\R^4)^n\to\C$ for $\zeta\in\phases\sqcup\{\mathrm{mix}\}$;
\item[H6.] functions $\mathsf F_n:\mathbb F^n\times\C^n\times(\R^3)^n\to\C$;
\item[H7.] finite sets $\channels{q}$, fields $O_{q,\lambda}\in F_q$ and functions $c_{q,\lambda}(\boldsymbol A;x,o)$ for every finite index set $I_{\rm cl}$, $\boldsymbol A\in\mathbb F^{I_{\rm cl}}$, $x\in(\R^4)^{I_{\rm cl}}$ and $o\in\R^4$;
\item[H8.] a field $\Theta\in F_4$.
\end{description}
\end{definition}

\begin{definition}[Reconstruction properties]\label{def:reconstruction}
A hierarchy satisfies the reconstruction properties when all the following hold:

\begin{enumerate}[label=\textup{R\arabic*.},leftmargin=3em,nosep]
\item Its relativistic datum satisfies W0–W3.
\item For every nonzero $\ell_{\rm lab}\in\mathcal L$, $\Phi(\ell_{\rm lab},f)\ne0$ for some $f\in\mathcal S$.
\item Each $\iota_d$ is injective.
\item $\Phi(\iota(A),f)\ne0$ for some $A\in\mathbb F$, $f\in\mathcal S$.
\item Every $\ell_{\rm lab}\in\mathcal L$ equals $\iota_dA$ for some $d,A$.
\item $\pi_{\rm quo}$ is surjective onto $\mathbb F$.
\item $\nu(\phases)=1$.
\item $\nu(O)>0$ for every nonempty open $O\subset\phases$.
\item For every state $\zeta$, simultaneous permutation of field and position slots preserves $S_n^\zeta$ on $\mathrm{Conf}_n$.
\item For every state $\zeta$ and field list $\boldsymbol A$, $S_n^\zeta(\boldsymbol A;\cdot)$ is smooth on $\mathrm{Conf}_n$.
\item For every fixed $n$, $\boldsymbol A\in\mathbb F^n$ and $x\in\Conf{n}$, $\xi\mapsto S_n^\xi(\boldsymbol A;x)$ is continuous.
\item $\int_\phases S_n^\xi(\boldsymbol A;x)\,d\nu(\xi)=S_n^{\mathrm{mix}}(\boldsymbol A;x)$ for all $n$, $\boldsymbol A\in\mathbb F^n$ and $x\in\Conf{n}$.
\item At each fixed real spatial configuration, $\boldsymbol z_{\rm time}\mapsto\mathsf F_n(\boldsymbol A;\boldsymbol z_{\rm time},\vec x)$ is holomorphic on $\mathcal Z_n=\{\boldsymbol z_{\rm time}\in\C^n:\operatorname{Im}(\boldsymbol z_{\rm time})_1<\cdots<\operatorname{Im}(\boldsymbol z_{\rm time})_n\}$.
\item If $t(x_1)<\cdots<t(x_n)$, then $\mathsf F_n(\boldsymbol A;(it(x_i))_i,(\vec x_i)_i)=S_n^{\xi_*}(\boldsymbol A;x)$.
\item For every $n$, $\boldsymbol A\in\mathbb F^n$, real $h_1<\cdots<h_n$ and $f_i\in\mathcal S$, the limit below holds.
\end{enumerate}

The limit in R15 is
\begin{equation}\label{eq:reconstruction-boundary}
\lim_{r\downarrow0}\int\mathsf F_n(\boldsymbol A;(t(x_i)+irh_i)_i,(\vec x_i)_i)\prod_if_i(x_i)\,dx
=\langle\Omega,\Phi(\iota(A_1),f_1)\cdots\Phi(\iota(A_n),f_n)\Omega\rangle\,.
\end{equation}
For this clause alone, use the convention that the integral of a nonintegrable function is zero; no separate integrability assertion at positive height is included.
\end{definition}

\paragraph{Products and comparison fields.} The OPE replaces a cluster by single fields, independently of spectators and state.

\begin{definition}[OPE geometry and remainders]\label{def:ope}
For finite $I_{\rm cl},I_{\rm sp}$, $\boldsymbol A\in\mathbb F^{I_{\rm cl}}$, $\boldsymbol B\in\mathbb F^{I_{\rm sp}}$, write $|\boldsymbol A|_{\deg}=\sum_{i\in I_{\rm cl}}\deg A_i$ and concatenate the two lists in a fixed order in $S^\zeta(\boldsymbol A,x;\boldsymbol B,y)$. For $r>0$, $M_{\rm geo}\ge0$, $\eta>0$, a cluster $(x,o)$ satisfies $|x_i-o|\le M_{\rm geo}r$ and $|x_i-x_{i'}|\ge\eta r$ for $i\ne i'$. A spectator configuration is $\delta_{\rm sp}$-clear if $|o-y_j|\ge\delta_{\rm sp}$ and $|y_j-y_{j'}|\ge\delta_{\rm sp}$ for $j\ne j'$. Define
\begin{equation}\label{eq:ope-remainder}
\mathcal R_D^\zeta(\boldsymbol A,x;\boldsymbol B,y;o)=S^\zeta(\boldsymbol A,x;\boldsymbol B,y)-\sum_{q=0}^D\sum_{\lambda\in\channels{q}}c_{q,\lambda}(\boldsymbol A;x,o)S^\zeta(O_{q,\lambda}(o),\boldsymbol B(y))\,.
\end{equation}
The OPE property is precisely:

\begin{description}[font=\normalfont,leftmargin=3.8em,labelindent=0pt,nosep]
\item[O1.] for every $I_{\rm cl},\boldsymbol A,M_{\rm geo},\eta,\varepsilon,q,\lambda$ with $M_{\rm geo}\ge0$, $\eta,\varepsilon>0$, there is $r_0>0$ such that $|c_{q,\lambda}(\boldsymbol A;x,o)|\le r^{q-|\boldsymbol A|_{\deg}-\varepsilon}$ for every $0<r<r_0$ and every such cluster;
\item[O2.] for every $I_{\rm cl},I_{\rm sp},\boldsymbol A,\boldsymbol B,M_{\rm geo},\eta,\delta_{\rm sp},\varepsilon,D$ with $M_{\rm geo}\ge0$, $\eta,\delta_{\rm sp},\varepsilon>0$, there is $r_1>0$ such that $|\mathcal R_D^\zeta|\le r^{D+1-|\boldsymbol A|_{\deg}-\varepsilon}$ for all $0<r<r_1$, every such cluster, every $\delta_{\rm sp}$-clear spectator configuration and every state $\zeta$.
\end{description}
\end{definition}

The coefficients in H7 do not depend on spectators, state or truncation.

\begin{definition}[Trace test and finite Wick profiles]\label{def:trace-test}
A mode is $(m_{\rm sh},j_h,k_E,h_{\rm mom})$ with $m_{\rm sh}\ge0$, $j_h,k_E\in\N$ and continuous real $h_{\rm mom}$ on $\R^3$ satisfying $h_{\rm mom}(\lambda_{\rm dil}\mathbf k)=\lambda_{\rm dil}^{j_h}h_{\rm mom}(\mathbf k)$ for $\lambda_{\rm dil}>0$. Put
\begin{equation}\label{eq:mode-kernel}
K_{m_{\rm sh},k_E,h_{\rm mom}}(r)=\int_{\R^3}\frac{\omega_{m_{\rm sh}}(\mathbf k)^{k_E}h_{\rm mom}(\mathbf k)}{2\omega_{m_{\rm sh}}(\mathbf k)}e^{-r\omega_{m_{\rm sh}}(\mathbf k)}\,d\mathbf k,\qquad \omega_{m_{\rm sh}}(\mathbf k)=\sqrt{m_{\rm sh}^2+|\mathbf k|^2}\,.
\end{equation}
A finite Wick profile is any finite complex linear combination of finite products of these kernels. A mass-shell integrand is assigned value zero at $\mathbf k=0$ if its denominator vanishes; this changes no integral. Define $\tracetwopoint(r)=S_2^{\xi_*}(\Theta,\Theta;-re_4/2,re_4/2)$. The trace test is:

\begin{description}[font=\normalfont,leftmargin=3.8em,labelindent=0pt,nosep]
\item[T1.] $\langle\Omega,\Phi(\iota_4\Theta,f)\Omega\rangle=0$ for every $f\in\mathcal S$;
\item[T2.] $\traceamplitude=3\dimG/\pi^4>0$;
\item[T3.] $r^8\log^2(1/r)\tracetwopoint(r)\to \traceamplitude$ as $r\downarrow0$, in $\C$;
\item[T4.] for every finite Wick profile $W_{\rm Wick}$ and every $\delta>0$, some $0<r<\delta$ has $\tracetwopoint(r)\ne W_{\rm Wick}(r)$.
\end{description}
\end{definition}

\begin{definition}[Polynomial mass-shell comparisons]\label{def:polynomial-comparisons}
For $m_{\rm sh}\ge0$ and a complex polynomial $P_{\rm mom}$ on $\R^4$, define the elementary momentum functional and its density by
\begin{equation}\label{eq:momentum-functional}
L_{m_{\rm sh},P_{\rm mom}}(\varphi)=\int_{\R^3}v_{m_{\rm sh},P_{\rm mom}}(\mathbf k)\varphi(\mathbf k,\omega_{m_{\rm sh}}(\mathbf k))\,d\mathbf k,\qquad v_{m_{\rm sh},P_{\rm mom}}(\mathbf k)=\frac{P_{\rm mom}(\mathbf k,\omega_{m_{\rm sh}}(\mathbf k))}{2\omega_{m_{\rm sh}}(\mathbf k)}\,.
\end{equation}
Finite polynomial comparison data on a set of generator labels assign a finite list $\mathcal M_{\rm pair}(\mathsf x,\mathsf y)$ of mass--polynomial pairs to each ordered pair of labels, and a complex number $c_{\rm pair}(u,u',\sigma)$ to each pair of words $u=(\mathsf x_1,\ldots,\mathsf x_{n_{\rm word}})$, $u'=(\mathsf y_1,\ldots,\mathsf y_{n_{\rm word}})$ and each permutation $\sigma$. Put
\begin{multline}\label{eq:polynomial-pairing}
\Pi^{\rm pair}(u,u';\varphi)=\sum_\sigma c_{\rm pair}(u,u',\sigma)\sum_{\substack{(m_i,P_i)\in\mathcal M_{\rm pair}(\mathsf x_i,\mathsf y_{\sigma i})\\1\le i\le n_{\rm word}}}\int_{(\R^3)^{n_{\rm word}}}\prod_i v_{m_i,P_i}(\mathbf k_i)\\
{}\cdot\varphi\Big(\sum_i(\mathbf k_i,\omega_{m_i}(\mathbf k_i))\Big)\,d\boldsymbol k\,.
\end{multline}
For words of different lengths put $\Pi^{\rm pair}=0$. A polynomial field $\mathsf Q$ is a finite list of words $u_i$ with complex amplitudes $\alpha_i^{\rm pol}$. Empty words are allowed. For two polynomial fields put
\begin{equation}\label{eq:polynomial-correlation}
W^{\rm pol}(\mathsf Q,\mathsf Q';\varphi)=\sum_{i,j}\alpha_i^{\rm pol}{\alpha'_j}^{\rm pol}\Pi^{\rm pair}(u_i,u'_j;\varphi)\,.
\end{equation}
With $\fourierplus{f}(\mathbf p)=\int e^{i\mathbf p\cdot x}f(x)\,dx$, the relative transform is
\begin{equation}\label{eq:relative-fourier}
(\wavetransform\psi)(\mathbf p)=\int e^{i\mathbf p\cdot(y-x)}\psi(x,y)\,dx\,dy,\qquad \wavetransform(\bar f\otimes h)=\overline{\fourierplus{f}}\fourierplus{h}\,.
\end{equation}
\end{definition}

\begin{definition}[Maxwell Fock comparisons]\label{def:fock-comparisons}
Put $d\sigma(\mathbf k)=d^3\mathbf k/(2|\mathbf k|)$ and $\mathsf P^\perp_{ij}(\mathbf k)=\delta_{ij}-\mathbf k_i\mathbf k_j/|\mathbf k|^2$, assigning the density and projector value zero at $0$. For $n_{\rm col}$ colours put
\begin{equation}\label{eq:fock-space}
\mathcal K_{n_{\rm col}}=L^2(\R^3,d\sigma;\C^{n_{\rm col}}\otimes\C^3),\qquad \mathcal F_{n_{\rm col}}=\bigoplus_{k\ge0}\mathcal K_{n_{\rm col}}^{\otimes k}\,.
\end{equation}
Let $\mathcal F_{n_{\rm col}}^{\rm phys}$ be the symmetric subspace transverse in every momentum slot. Translations multiply a $k$-particle wave by $e^{iz_{\rm tr}\cdot\sum_j(\mathbf k_j,|\mathbf k_j|)}$. The vacuum $\Omega_0$ lies in sector zero. For a list $\boldsymbol h=(h_{\alpha,i})$ of Schwartz tests put
\begin{equation}\label{eq:coherent-vector}
v_{\boldsymbol h}(\alpha,i;\mathbf k)=|\mathbf k|\sum_j\mathsf P^\perp_{ij}(\mathbf k)\widehat{h}_{\alpha,j,+}(\mathbf k,|\mathbf k|),\qquad \mathcal E(\boldsymbol h)=e^{-\|v_{\boldsymbol h}\|^2/2}\bigoplus_{k\ge0}\frac{v_{\boldsymbol h}^{\otimes k}}{\sqrt{k!}}\,.
\end{equation}
A Fock comparison is a complex linear isometry $\mathsf I_{\rm F}:\mathcal H\to\mathcal F_{n_{\rm col}}$ with image in $\mathcal F_{n_{\rm col}}^{\rm phys}$, $\mathsf I_{\rm F}\Omega=\Omega_0$, and intertwining translations. Write $\psi_{\mathsf I_{\rm F}}(f)=\mathsf I_{\rm F}\Phi(\iota_4\Theta,f)\Omega$. For the averaged comparison, a probability space $(Z,\lambda_Z)$ and maps $P_z^{\rm av}:\mathcal F_{n_{\rm col}}\to\mathcal F_{n_{\rm col}}$ preserve inner products on pairs of physical vectors and fix $\mathsf I_{\rm F}\mathcal H$ pointwise. Linearity of $P_z^{\rm av}$ is not required. The scalar integrals use the convention of Definition~\ref{def:reconstruction}.
\end{definition}

\paragraph{Nested products and contacts.} Successive OPE expansions use relative scales.

\begin{definition}[Nested clusters and their rates]\label{def:nested-clusters}
Let a finite rooted tree have leaves carrying fields $A_i$ and an internal root $o_{\rm root}$. Assign a position $z_v$ to every vertex and a positive scale $\varrho_v$ to every internal vertex, with $\varrho_{o_{\rm root}}=r$ and $\varrho_v\le\varrho_{\operatorname{par}(v)}$. At each internal $v$, every descendant position lies within $M_{\rm geo}\varrho_v$ of $z_v$, and positions belonging to different immediate child subtrees have distance at least $\eta\varrho_v$. These conditions apply recursively. Spectator clearance $\delta_{\rm sp}$ means distance at least $\delta_{\rm sp}$ from every tree position and between distinct spectators. Put $W_v=\sum_{i\text{ leaf below }v}\deg A_i$ and $q_v=\varrho_v/\varrho_{\operatorname{par}(v)}$. For truncations $\boldsymbol D=(D_v)$ define
\begin{equation}\label{eq:nested-rates}
\mathcal Q_{\varepsilon,q}(\varrho)=r^{q-W_{o_{\rm root}}-\varepsilon}\prod_{\substack{v\ne o_{\rm root}\\v\text{ internal}}}q_v^{-W_v-\varepsilon},\qquad \mathsf e_{\boldsymbol D}(\varrho)=\sum_{\substack{v\ne o_{\rm root}\\v\text{ internal}}}q_v^{D_v+1}\,.
\end{equation}
The direct coefficient $C^{\rm dir}_{q,\lambda}$ is $c_{q,\lambda}$ of all original leaves at $z_{o_{\rm root}}$. Define $C^{\rm it,\boldsymbol D}_{q,\lambda}$ recursively. A leaf $i$ carries $(A_i,1)$. A proper internal vertex $v$ with children $c_1,\ldots,c_s$ carries $(O_{q',\lambda'},\kappa_v(q',\lambda'))$ for $q'\le D_v$, $\lambda'\in\channels{q'}$, where
\begin{equation}\label{eq:iterated-coefficient}
\kappa_v(q',\lambda')=\sum_{(B_j,\kappa_j)_j}\Bigl(\prod_{j=1}^s\kappa_j\Bigr)c_{q',\lambda'}(B_1,\ldots,B_s;z_{c_1},\ldots,z_{c_s},z_v)\,.
\end{equation}
The sum chooses one pair $(B_j,\kappa_j)$ from each child. The same formula at the root gives $C^{\rm it,\boldsymbol D}_{q,\lambda}=\kappa_{o_{\rm root}}(q,\lambda)$, without using the root truncation. A tested root remainder in state $\zeta$ is $\mathcal R_q^\zeta$ for the original leaves at $z_{o_{\rm root}}$.
\end{definition}

\begin{definition}[Contact distributions]\label{def:contacts}
A distribution on $\R^{4n}$ is a complex linear continuous functional on $C_c^\infty(\R^{4n})$. Locally finite order means that each compact $K$ admits $q_K\in\N$, $C_K<\infty$ with
\begin{equation}\label{eq:contact-order}
|\mathfrak D(f)|\le C_K\max_{|\alpha|\le q_K}\|\partial^\alpha f\|_\infty
\end{equation}
for tests supported in $K$. A contact extension $\mathfrak D_{\boldsymbol A}^\zeta$ of $S^\zeta(\boldsymbol A;\cdot)$ is such a distribution equal to $\int fS^\zeta(\boldsymbol A;\cdot)$ on separated tests. The subtraction prescription is defined in Section~\ref{sec:partII-J}, using the field projections of the OPE.
\end{definition}

Continuity gives locally finite order for every distribution; Proposition~\ref{prop:P13} records the dependence of $q_K,C_K$.
